\section*{第 \ref{ch:g2:balance-equilibrium} 章——平衡}

\begin{solution}{exo:g2:balance-equilibrium:1}
地球把台灯向下拉，桌面向上推它，
两力大小恰好相等。两力互相\omterm{def:g2:balance-equilibrium:equilibrium}{平衡}——即处于\omterm{def:g2:balance-equilibrium:equilibrium}{平衡}状态，所以台灯保持不动。
\end{solution}

\begin{solution}{exo:g2:balance-equilibrium:2}
两队朝相反方向用力且力气相等时，拉力相互抵消。
一旦一队的力气大于另一队，
\omterm{def:g2:balance-equilibrium:equilibrium}{平衡}就被打破。
\end{solution}

\begin{solution}{exo:g2:balance-equilibrium:3}
跷跷板保持水平——两边重量相同，距离也相同。当其中一个双胞胎
往外挪远，那一侧就会\omterm{def:g1:floating-and-sinking:words}{下沉}：距离同样起作用。
\end{solution}

\begin{solution}{exo:g2:balance-equilibrium:4}
靠近支点坐，莉娜则远远地坐在她那一侧。规律：
较重的人靠近中央，可以与较轻的人远坐时保持\omterm{def:g2:balance-equilibrium:equilibrium}{平衡}，
重量与距离两者都起作用。
\end{solution}

\begin{solution}{exo:g2:balance-equilibrium:5}
尺子的\omterm{def:g2:balance-equilibrium:point}{平衡点}在中间；锤子的\omterm{def:g2:balance-equilibrium:point}{平衡点}靠近
锤头一端。验证方法：用一根手指各顶住它们，
找到使其保持水平的那个点。
\end{solution}

\begin{solution}{exo:g2:balance-equilibrium:6}
\omterm{def:g2:balance-equilibrium:point}{平衡点}就在靠近刷子的一端，也就是较重的那一端。
\omterm{def:g2:balance-equilibrium:point}{平衡点}总是偏向物体较重的一侧。
\end{solution}

\begin{solution}{exo:g2:balance-equilibrium:7}
宽底的金字塔可以倾斜很大角度，直到重力
的作用线超出宽阔的底座。细窄的塔稍一倾斜就会倒下。
\end{solution}

\begin{solution}{exo:g2:balance-equilibrium:8}
\omterm{def:g2:balance-equilibrium:equilibrium}{平衡}杆能减缓每一次晃动，让走钢丝的人有时间
重新回到绳子上方。走在路缘上时伸开双臂也是同样道理：
晃动变慢，更容易把身体稳住。
\end{solution}

\begin{solution}{exo:g2:balance-equilibrium:9}
若只用右脚站立，身体必须向右倾，
使重力落在那只脚的正上方，正如塔使重力保持在底座之上。
墙恰好就是你需要靠的地方：
它阻止了这一动作，你的重心仍悬在
抬起那只脚的空位上方，于是你会倾倒。
\end{solution}

\begin{solution}{exo:g2:balance-equilibrium:10}
起重机就像一个巨大的跷跷板，塔身是支点。第一方：
吊起的重物远远地位于长前臂上。第二方：
混凝土块——重得多，所以可以靠近支点，坐在较短的后臂上。
重而近与轻而远相互\omterm{def:g2:balance-equilibrium:equilibrium}{平衡}，正是跷跷板的规则；
后臂之所以短，是因为其上的“坐者”最重。
\end{solution}
